\section{Ameaças à Validade e Mitigações}

{\footnotesize
\arrayrulecolor{BorderGray}
\renewcommand{\arraystretch}{1.35}
\setlength{\tabcolsep}{5pt}
\begin{longtable}{|L{7.65cm}|L{8.95cm}|}
\hline
\rowcolor{HeaderBg}
\tblhead{Ameaça} & \tblhead{Mitigação} \\
\hline
\endfirsthead
\hline
\rowcolor{HeaderBg}
\tblhead{Ameaça} & \tblhead{Mitigação} \\
\hline
\endhead
\hline
\endfoot
\textbf{Não resposta:} quem tem sobrecarga ou desinteresse em IA não responde; na bola de neve a taxa de resposta é incalculável. & Questionário curto ($\sim 4$ min); comparação entre respostas iniciais e tardias. \\ \hline
\textbf{Autosseleção:} sobre-representação de entusiastas de IA. & Chamada neutra (\enquote{estudo sobre práticas reais e impactos de ferramentas}); ITEM\_07 mostra a intensidade de uso na amostra. \\ \hline
\textbf{Desejabilidade social:} constrangimento ao admitir aceitar código sem validar. & Anonimato total e alternativas neutras (reduz, sem eliminar). \\ \hline
\textbf{Validade de construto:} produtividade e manutenibilidade medidas por poucos itens; M\_REV tem 2 itens. & Construtos multi-item, sem itens compostos; alfa e $\rho$ (que medem confiabilidade, não validade); piloto com 4 pessoas é limitação declarada. \\ \hline
\textbf{Autorrelato e método comum:} tudo é percepção do mesmo respondente, o que pode inflar correlações (QP4). & Item invertido, ordem aleatória, análises de sensibilidade, linguagem restrita a \enquote{percepção}. \\ \hline
\textbf{Recordação:} comparação com um \enquote{antes} lembrado (ITEM\_11, ITEM\_14). & Limitação reconhecida, sem mitigação pelo desenho. \\ \hline
\textbf{Validade externa:} amostra não probabilística (LinkedIn/Telegram/WhatsApp), só Brasil. & Sem generalização estatística; perfil da amostra descrito (M\_SEN, M\_TOOL, M\_USO). \\ \hline
\textbf{Conclusão:} baixo poder ($n \approx 60$, grupos desbalanceados) e múltiplos testes. & Apenas 3 testes confirmatórios (QP3 com $\alpha = 0{,}025$ cada; QP4 com $\alpha = 0{,}05$); regra de fusão de grupos; tamanhos de efeito com ICs. \\ \hline
\end{longtable}
\arrayrulecolor{black}
}
